\section{Unified Mobile Governed Communication}
\label{sec:pilot-unified-mobile-communication}

The unified Pilot mobile application now projects the existing governed communication authority
onto a dedicated private Messages surface.  The phone is a possession-bound review and approval
surface; it is not an independent messaging server.  Contacts, drafts, provider bindings and
receipts remain on the person's mother brain, while shared television speech may prepare a request
but cannot choose a private recipient, sending account or final message by itself.

\subsection{User Experience and Action States}

The Personal space includes an Email and Messages shortcut.  From this surface the person can
select one of their private contacts, choose an available email, WhatsApp or Pilot-to-Pilot route,
enter a subject where applicable, and prepare the exact content.  Preparation creates no external
effect.  The draft is then displayed with its resolved contact name, exact private destination,
subject, body and current state.

The delivery lifecycle is deliberately explicit:
\begin{enumerate}
  \item \texttt{awaiting\_private\_approval}: the recipient and content are visible for review;
  \item \texttt{guardian\_approval\_required}: a child or teenager cannot approve an external route alone;
  \item \texttt{approved\_pending\_adapter}: the exact hash was approved, but nothing has been sent;
  \item \texttt{executing}: a single authorized provider attempt may be in flight;
  \item \texttt{delivery\_unknown\_reconcile\_required}: Pilot cannot safely claim either success or failure;
  \item \texttt{delivered}: the provider returned a verified message identifier and Pilot recorded a receipt; and
  \item \texttt{rejected}: the person declined the draft and no delivery occurred.
\end{enumerate}

Approval and execution remain separate taps.  If no authorized sending account or platform adapter
is connected, Pilot keeps the approved draft intact and reports that delivery is unavailable.  It
does not open an uncontrolled browser, invent success or downgrade the action into an unverified
claim.

\subsection{Mother-Brain Architecture}

The mobile surface calls a provider-independent communication projection through signed HTTPS
routes.  That projection reuses \texttt{GovernedCommunication}, \texttt{PilotInbox}, the private
contact authority, production identity records and persona-specific Google OAuth metadata.  The
mother brain performs recipient resolution, immutable-scope hashing, policy checks, adapter
selection and receipt normalization.  Provider credentials and OAuth vault references are never
returned to the phone or television.

The current provider boundary supports a Gmail sending adapter that requires the person's own
authorized Google binding.  WhatsApp and equivalent messaging delivery remain disabled until an
officially permitted provider adapter is configured.  Pilot-to-Pilot conversation transport remains
separate from public-provider messaging and continues to use the encrypted Unified Inbox route.

\subsection{Possession-Bound Capabilities}

Communication access is divided into narrow lease scopes:
\begin{itemize}
  \item \texttt{communication.read} for private drafts, summaries, follow-ups and receipts;
  \item \texttt{communication.draft} for preparing exact content without sending it;
  \item \texttt{communication.approve} for accepting or rejecting the unchanged content hash;
  \item \texttt{communication.execute} for one approved provider attempt;
  \item \texttt{communication.convert} for turning an authorized received summary into a task or calendar proposal;
  \item \texttt{communication.followup} for a conditional no-reply reminder; and
  \item \texttt{communication.call} for preparing a telephone handoff that still requires a deliberate phone tap.
\end{itemize}

Every mobile operation is signed with the phone's non-exportable Ed25519 key and checked against
the active certificate, persona and short-lived capability lease.  Persona substitution fails
closed.  Draft and decision retries use idempotency keys bound to canonical request hashes; an
identical retry returns the original object, while reuse with changed content or a changed decision
is rejected.

\subsection{Verified Delivery and Minimal Receipts}

Execution is permitted only when the same persona approved the stored content hash.  A successful
adapter must return a verified provider message identifier.  Pilot then stores a normalized receipt
containing the draft identifier, persona, channel, provider message and thread identifiers, a hash
of the recipient, and the verified delivery time.  Raw provider responses are not retained.  A
repeat execution returns the original receipt rather than sending the message twice.

If the adapter raises an error after an attempt begins, Pilot records an unknown-effect state and
requires reconciliation.  This prevents an uncertain timeout from causing an automatic duplicate
message.  The interface never labels a draft as delivered merely because it left the phone or was
accepted by the mother brain.

\subsection{Received Summaries, Follow-Through and Calls}

Authorized provider ingestion stores sender, freshness and source provenance together with a
bounded AION Native summary and a hash of the original body.  The raw body is not retained in this
communication record.  From the private phone, the person may convert a received summary into a
local task or prepare a calendar item; calendar execution still uses its own exact approval gate.

After a verified sent message, the person may schedule a reminder conditioned on the absence of a
verified reply.  Pilot also supports call preparation using a privately stored telephone route.
Preparation never places the call: the phone presents a separate \emph{Tap to call} action so the
operating system and user remain the final call authority.

\subsection{Verification and Remaining External Qualification}

Automated tests cover private snapshots, retry-safe drafting, changed-content rejection, exact
approval, missing-adapter failure, verified provider receipts, repeat-execution deduplication,
conditional follow-ups, received-message task conversion, provenance retention, raw-body exclusion
and explicit call placement.  Surface checks cover the Messages shortcut, exact-content review,
approval, execution and verified receipt views.  At coded closure of \texttt{MOB-0506}, all 460
AION Fabric tests passed and the optimized Next.js application compiled and statically generated
the unified mobile route.

No real message, email or call was made during this implementation.  Production qualification still
requires a user-authorized Google account and a real Gmail receipt, plus separately authorized and
officially supported messaging providers where available.  Those external qualifications remain
visible in the master checklist and cannot be closed by a simulated adapter.
